\section{The Postcard and the Measuring Reed}

There is a moment near the end of the Apocalypse that nobody paints. An angel arrives with a measuring rod---``a measure of a reed of gold, to measure the city and the gates thereof and the wall''---and takes the dimensions of heaven. The result is reported with the flatness of a surveyor's docket: ``And the city lieth in a four-square: and the length thereof is as great as the breadth. And he measured the city with the golden reed for twelve thousand furlongs: and the length and the height and the breadth thereof are equal.''\endnote{Apocalypse 21:15--16, Douay--Rheims (Challoner revision), read 2026-07-26 in the source library's tracked verse text of this edition at the registered artifact for the Apocalypse (lines 373--374 of the per-verse file). Scripture throughout this article is quoted from that tracked witness and cited in the Vulgate chapter-and-verse numbering the edition preserves; where that numbering differs from a modern Bible's, the note says so. The conversion of twelve thousand \latin{stadia} to a modern distance is the article's own arithmetic at the conventional 185-metre \latin{stadion}, and is offered only to make the figure imaginable; the text gives no metric equivalent and the number is plainly symbolic (twelve, the tribes and the apostles, multiplied by a thousand).} Height equal to breadth: the holy city is a cube about fourteen hundred miles on a side, with walls of jasper and twelve gates each cut from a single pearl, and no temple in it, and no sun.

Set that beside the picture most people actually carry. In the picture there are clouds, and the clouds are load-bearing---souls stand on them. The souls have wings, or are about to be issued some. There is a gate with a saint behind a lectern, checking names. There are harps, one to a customer. There is a reunion: the people you loved, restored to you at their best ages, in a place that is recognisably a very good version of here. The picture is not wicked, and it is not empty. But almost none of it is in the Apocalypse, and the parts that are in the Apocalypse are not in it the way the picture assumes. The gates of pearl are there, but they stand open and are never shut, and nobody is stationed at them with a list.\endnote{Apoc 21:21 (``the twelve gates are twelve pearls, one to each: and every several gate was of one several pearl'') and 21:25 (``And the gates thereof shall not be shut by day: for there shall be no night there''); 21:12 places twelve angels in the gates, with the names of the twelve tribes written on them. The figure of a single apostle admitting or refusing entrants at a gate is not in the text; the keys given to Peter in Matt 16:19 are keys of the kingdom given on earth, and the Apocalypse does not station him at the gates. All read 2026-07-26 in the tracked verse text.} Harps are there, but the first hands holding them are the four living creatures and the twenty-four ancients, and when the redeemed are given harps they are given ``the harps of God'' while standing on a sea of glass mingled with fire, which is not what the picture means by a harp.\endnote{Apoc 5:8; 14:2; 15:2, tracked verse text, read 2026-07-26. On the wings, see the bounded search reported in the section on Scripture and the postcard below.} Wings are not there at all---not once, in the whole New Testament of this edition, given to a human being living or risen.

That gap is the subject of this article, and it runs in both directions, which is what makes it interesting. The devotional picture supplies a great deal that revelation never gave. And revelation, together with the Church's definitions and the tradition's hardest thinking, contains a great deal that the picture has quietly dropped---beginning with the fact that the central Catholic claim about heaven is not about a place at all, and is far stranger, far more specific, and far harder to believe than clouds.

\subsection{Where this volume stands}

This is the last of four studies on the four last things. Its companions took the three subjects that happen \emph{to} a person: the death that no one can delegate; the judgment in which a life is finally told the truth about itself; and the refusal that can become final. Each of those is, in a sense, an event. Heaven is not an event. It is a state, and the tradition's word for it---\latin{beatitudo}, beatitude, blessedness---names not a location or a duration but a condition of a creature that has got what it is for.\endnote{The other three volumes of this series are independent works and are not summarised or relied on here; where their subjects touch this one---the particular judgment that precedes the state described in \emph{Benedictus Deus}, the purification that the same constitution presupposes, the possibility of definitive refusal---this article states only what its own sources establish and leaves those subjects to the volumes that own them. \latin{Beatitudo} is rendered ``beatitude'' or ``happiness'' throughout; it translates the complete and final good of a rational creature, not a mood or a pleasant feeling, and the English word ``happiness'' in quotations from the English Dominican \emph{Summa} carries that heavier sense.} So the closing volume of a series about last things cannot simply add a fourth fact to three facts. It has to say what the whole sequence was for---and it has to do that without pretending to know more than it does.

\subsection{Three registers, and why keeping them apart is the whole job}

Almost everything a Catholic reader has ever heard about heaven belongs to one of three registers, and the registers have wildly different weights.

The first is what the Church has \emph{defined}: a small number of propositions, promulgated in acts whose texts can be read, whose wording was fought over, and which bind. The core of it fits comfortably on one page, and it will be quoted here in full.

The second is what the tradition has \emph{argued}: the immense and often magnificent structure of theological reasoning that stands on the defined core---Augustine on the risen body, Aquinas on the light of glory, Gregory Palamas on the uncreated energies, the schoolmen on the endowments of the glorified flesh. This material is not idle. Some of it is close to unanimous; some of it is the common teaching of theologians and can be denied only rashly; some of it is one school's opinion, contradicted by another school in good standing. It is all reasoning, and reasoning can be wrong.

The third is what devotion has \emph{imagined}: the clouds, the wings, the gate-keeper, the reunion scene. This register is not simply an error to be swept away---images are how most people hold onto anything, and the Church herself teaches in images because Scripture does.\endnote{\emph{Catechism of the Catholic Church} 1027: ``This mystery of blessed communion with God and all who are in Christ is beyond all understanding and description. Scripture speaks of it in images: life, light, peace, wedding feast, wine of the kingdom, the Father's house, the heavenly Jerusalem, paradise\ldots'' Official Vatican English web text, read 2026-07-26; the Catechism's own delivery of the passage is quoted, and its typography (a spaced hyphen for the dash) is not reproduced. Quotations from the Catechism throughout are brief excerpts of a copyrighted official text, used with attribution.} The trouble begins only when an image is cashed as a description---when ``the Father's house'' stops being a promise of welcome and becomes a floor plan.

The discipline this article tries to keep is simple to state and unpleasant to apply: at every consequential claim, say which register it belongs to. The unpleasantness is that the honest answer shrinks the picture. A reader who wants heaven described will get less here than a devotional book would give. What is offered instead is a smaller thing that can bear weight: a certain core, marked as certain; a body of serious reasoning, marked as reasoning, with the disagreements left standing where the tradition left them standing; and a candid account of how much of the received picture is neither.

That last figure is worth stating at the outset, because it governs everything that follows. On the evidence assembled here, the great majority of what an ordinary Catholic ``knows'' about heaven---what the risen body will look like, how old it will be, what the blessed will do, whether they see one another, what the new earth will be like, how any of it feels---is theological speculation, and a good deal of it is speculation that its own authors explicitly marked as such and that later readers un-marked.\endnote{This is the article's own summary judgment, a project synthesis rather than a received teaching; the evidence for it is the whole of the sections that follow, and it is stated again in bounded form, with its counterargument, in the synopsis at the end. The most striking single piece of evidence is Augustine's own warning in the two chapters of \emph{De civitate Dei} on which most of the Western picture of the risen body ultimately rests: ``here we speak of a matter of which we have no experience, and upon which the authority of Scripture does not definitely pronounce'' (XXII.29) and ``what this spiritual body shall be, and how great its grace, I fear it were but rash to pronounce, seeing that we have as yet no experience of it'' (XXII.21). Both quoted in full below from the tracked 1871 Dods translation.} Augustine, who supplied more of that picture than anyone, stopped twice in the middle of supplying it to say that he was guessing. Aquinas, reading Augustine on whether risen eyes will see God, notes drily that Augustine ``speaks as one inquiring, and conditionally.''\endnote{Thomas Aquinas, \emph{Summa theologiae} I, q.~12, a.~3, reply to objection 2, English Dominican translation, read 2026-07-26 at New Advent. The New Advent presentation inserts spaces before punctuation where a word is hyperlinked; quotations here close those spaces and change nothing else. Latin checked at Corpus Thomisticum for the loci identified in the References.} The tentativeness was in the sources from the beginning. It fell out somewhere between the sources and the holy card.

\subsection{What is coming}

The argument runs in the order the evidence actually has. It begins with the definitions, because they are the fixed points and because the story of how one of them was made---a pope preaching one thing and his successor defining the opposite, two years after his death---is the best available lesson in what a definition is and is not. It then asks the question the definitions raise and do not answer: whether a creature can see God at all, and if so how, working the Thomistic answer directly in the \emph{Summa} rather than through summaries of it. It then puts the Christian East's very different answer at its full strength, in its own authors, because that answer is not a curiosity but a live rival that reads the same Scripture and the same Fathers. It then takes up the body that rises, the degrees of glory, and the company---three places where the tradition's confident detail and its actual warrant diverge most sharply. It closes by auditing the postcard against Scripture, and then by asking what the four last things add up to.

The governing claim, stated now so that the reader can watch it be tested rather than sprung at the end: what the Catholic Church has defined about heaven is a promise of an immediate, unmediated, personal seeing of God himself, given to creatures who cannot achieve it, differing in degree with charity, entailing the resurrection of the very bodies we now bear, and enjoyed in company---and that this defined core is both smaller and more extreme than the picture that has grown over it. Smaller, because nearly every furnishing has been added. More extreme, because the picture's clouds and reunions are, by comparison, modest requests.
